\chapter{Methodology}\label{ch:method}

\section{Generated Data}
The script \texttt{scripts/make\_data.py} creates 2,880 hourly rows for three example sites. The temperature and power columns contain 50 deliberately inserted missing values in total. The figure script fills those gaps by time-ordered linear interpolation. The dataset is synthetic; it is generated to exercise the plotting and typesetting workflow.

\section{Hinge Regression}
The example fits the hinge regression\index{hinge regression} model in \Cref{eq:thesis-hinge}:
\begin{equation}
  P_k=\beta_0+\beta_1\max(T_k-\tau,0)+\varepsilon_k,
  \label{eq:thesis-hinge}
\end{equation}
where $T_k$ is temperature, $P_k$ is power, and $\tau$ is the threshold. The script uses a grid search\index{grid search} over candidate thresholds. For each candidate, it solves for $\beta_0$ and $\beta_1$ by ordinary least squares\index{ordinary least squares} and keeps the candidate with the lowest residual sum of squares.

\begin{proposition}\label{prop:ols}
For a full-column-rank design matrix $\bm X$, the ordinary least-squares estimator
\begin{equation}
  \hat{\bm\beta}
  =\argmin_{\bm\beta}\lVert\bm y-\bm X\bm\beta\rVert_2^2
  =(\bm X^{\top}\bm X)^{-1}\bm X^{\top}\bm y
  \label{eq:thesis-ols}
\end{equation}
is unbiased under the zero-mean error assumption. If the errors have covariance $\sigma^2\bm I$, then
$\Cov(\hat{\bm\beta})=\sigma^2(\bm X^{\top}\bm X)^{-1}$.
\end{proposition}

\section{Control-Diagram Examples}
\Cref{fig:loop-thesis} shows the structure of a feedback loop. The PID equation below defines the controller signal in terms of the error:
\begin{equation}
  u(t)=K_p e(t)+K_i\!\int_0^t e(\tau)\,d\tau+K_d\dot e(t).
\end{equation}
The second-order step responses\index{step response} in \cref{fig:step-thesis} compare damping ratios of 0.2, 0.5, and 0.8. The lower-damped response has more overshoot. These plots illustrate control concepts; no controller is fitted to or tested on the generated sensor data.

\begin{figure}[htbp]
  \centering
  \resizebox{\textwidth}{!}{% Closed-loop feedback control system with PID controller, plant and sensor.
\begin{tikzpicture}[node distance=9mm and 9mm, auto]
  \node[coordinate] (in) {};
  \node[sum, minimum size=5mm, right=10mm of in] (cmp) {};
  \node[block, right=of cmp] (pid) {PID\\controller\\[-1pt]$C(s)$};
  \node[block, right=of pid] (act) {Actuator\\$A(s)$};
  \node[block, right=of act] (plant) {Plant\\$G(s)$};
  \node[coordinate, right=12mm of plant] (out) {};
  \node[block, below=10mm of act] (sens) {Sensor\\$H(s)$};
  \node[coordinate, above=9mm of plant] (dist) {};
  \draw[flow] (in) -- node[pos=0.3] {$r(t)$} (cmp);
  \draw[flow] (cmp) -- node {$e(t)$} (pid);
  \draw[flow] (pid) -- node {$u(t)$} (act);
  \draw[flow] (act) -- (plant);
  \draw[flow] (plant) -- node[pos=0.75] {$y(t)$} (out);
  \draw[flow] ($(plant.east)+(6mm,0)$) |- (sens.east);
  \draw[flow] (sens.west) -| node[pos=0.25, swap] {$y_m(t)$} (cmp.south);
  \draw[flow] (dist) -- (plant.north);
  \node[left=1mm of dist, font=\scriptsize] {disturbance $d(t)$};
  \node[font=\scriptsize] at ($(cmp.north west)+(-1mm,1.5mm)$) {$+$};
  \node[font=\scriptsize] at ($(cmp.south east)+(2.2mm,-1.5mm)$) {$-$};
  \node[above=1mm of pid, font=\scriptsize, text=black!70, align=center]
       {$u=K_p e+K_i\!\int e\,d\tau+K_d\dot e$};
\end{tikzpicture}
}
  \caption{Block diagram of a feedback loop with a PID controller.}
  \label{fig:loop-thesis}
\end{figure}

\begin{figure}[htbp]
  \centering
  % Step response of a second-order system for several damping ratios,
% plotted natively with PGFPlots (vector, fonts match the document).
\begin{tikzpicture}
  \begin{axis}[
      width=\linewidth, height=0.62\linewidth,
      xlabel={Time $t$ (s)}, ylabel={Response $y(t)$},
      xmin=0, xmax=12, ymin=0, ymax=1.7,
      grid=both, grid style={line width=.1pt, draw=gray!20}, major grid style={draw=gray!35},
      legend pos=south east, legend style={font=\scriptsize, draw=none, fill=white, fill opacity=0.85},
      every axis plot/.append style={line width=0.8pt}, samples=200, smooth, domain=0:12,
      tick label style={font=\footnotesize}, label style={font=\footnotesize}]
    \def\wn{1}
    \addplot[oiVerm]{1 - exp(-0.2*\wn*x)*(cos(deg(\wn*sqrt(1-0.04)*x)) + (0.2/sqrt(1-0.04))*sin(deg(\wn*sqrt(1-0.04)*x)))};
    \addlegendentry{$\zeta=0.2$}
    \addplot[oiBlue]{1 - exp(-0.5*\wn*x)*(cos(deg(\wn*sqrt(1-0.25)*x)) + (0.5/sqrt(1-0.25))*sin(deg(\wn*sqrt(1-0.25)*x)))};
    \addlegendentry{$\zeta=0.5$}
    \addplot[oiGreen]{1 - exp(-0.8*\wn*x)*(cos(deg(\wn*sqrt(1-0.64)*x)) + (0.8/sqrt(1-0.64))*sin(deg(\wn*sqrt(1-0.64)*x)))};
    \addlegendentry{$\zeta=0.8$}
    \addplot[black, dashed, domain=0:12, samples=2]{1};
    \addlegendentry{Setpoint}
  \end{axis}
\end{tikzpicture}

  \caption{Second-order unit-step responses for three damping ratios.}
  \label{fig:step-thesis}
\end{figure}

\section{Reproducible Figures}
The charts and summary table are generated by \texttt{scripts/generate\_figures.py}. The full sequence of commands is given in \cref{app:repro}.
